\section*{Data Pipelines}

Los datos se generan en bases de datos, aplicaciones, sensores, archivos,
servicios web y plataformas externas. Para que puedan utilizarse en análisis,
reportes, automatización o modelos de inteligencia artificial, es necesario
transportarlos y prepararlos de manera confiable. Los Data Pipelines permiten
automatizar este recorrido desde las fuentes hasta los sistemas que consumen la
información.

\section*{1. ¿Qué es un Data Pipeline?}

Un \textit{Data Pipeline} es un conjunto organizado y automatizado de procesos
que mueve datos desde una o varias fuentes hasta uno o varios destinos. Durante
el recorrido, la información puede ser validada, filtrada, limpiada,
transformada, enriquecida, combinada y almacenada.

Cada etapa recibe una entrada, realiza una operación y genera una salida que
sirve como entrada para la siguiente etapa. Dependiendo de las necesidades del
negocio, el flujo puede ejecutarse en horarios definidos, activarse cuando
ocurre un evento o procesar datos continuamente. Su finalidad es entregar
información correcta, oportuna y utilizable a sistemas analíticos u
operacionales.

\section*{2. ¿Cuáles son los tipos de Data Pipeline?}

Los Data Pipelines pueden clasificarse de distintas maneras.

\subsection*{Según la forma de procesamiento}

\begin{itemize}
    \item \textbf{Pipeline por lotes:} acumula datos durante un periodo y los
    procesa juntos en intervalos programados. Es apropiado para reportes
    históricos, cierres diarios, facturación y cargas nocturnas.

    \item \textbf{Pipeline de streaming:} procesa los datos de forma continua
    conforme se generan. Se utiliza cuando se requieren resultados con baja
    latencia, por ejemplo, para detectar fraude o monitorear sensores.

    \item \textbf{Pipeline de micro-lotes:} agrupa pequeñas cantidades de datos
    durante intervalos muy cortos. Representa un punto intermedio entre el
    procesamiento por lotes y el streaming continuo.

    \item \textbf{Pipeline híbrido:} combina procesos por lotes y streaming.
    Puede utilizar datos históricos para generar contexto y eventos recientes
    para producir respuestas inmediatas.
\end{itemize}

\subsection*{Según el orden de las transformaciones}

\begin{itemize}
    \item \textbf{ETL:} primero extrae, después transforma y finalmente carga
    los datos en el destino.

    \item \textbf{ELT:} extrae y carga rápidamente los datos originales; las
    transformaciones se ejecutan posteriormente dentro del sistema de destino.
\end{itemize}

\subsection*{Según el mecanismo de activación}

\begin{itemize}
    \item \textbf{Programado:} se ejecuta a una hora o con una frecuencia definida.
    \item \textbf{Controlado por eventos:} comienza cuando ocurre un evento determinado.
    \item \textbf{CDC:} utiliza captura de cambios para transportar únicamente las inserciones, actualizaciones y eliminaciones de una fuente.
\end{itemize}

\section*{3. ¿Cuáles son los elementos o componentes de un Data Pipeline?}

Los componentes principales son:

\begin{enumerate}
    \item \textbf{Fuentes de datos.} Son los sistemas donde se origina la
    información, como bases de datos, aplicaciones empresariales, API, archivos,
    dispositivos IoT, registros y servicios en la nube.

    \item \textbf{Conectores e ingestión.} Extraen o reciben los datos y los
    incorporan al pipeline. La ingestión puede realizarse por lotes, streaming o
    captura de cambios.

    \item \textbf{Transporte o mensajería.} Mueve eventos entre componentes y
    puede actuar como búfer para desacoplar productores y consumidores.

    \item \textbf{Procesamiento y transformación.} Limpia, valida, filtra,
    estandariza, agrega, combina y enriquece la información.

    \item \textbf{Almacenamiento temporal.} Conserva datos de manera
    intermedia para absorber variaciones de carga, permitir reintentos o
    mantener puntos de control.

    \item \textbf{Destino o \textit{sink}.} Es el sistema que recibe el
    resultado, por ejemplo, una base de datos, un Data Warehouse, un Data Lake,
    un Lakehouse, una aplicación o un modelo de aprendizaje automático.

    \item \textbf{Orquestación.} Define dependencias, orden de ejecución,
    horarios, reintentos y condiciones entre las tareas.

    \item \textbf{Calidad y validación.} Comprueba formatos, rangos, reglas de
    negocio, duplicados, valores faltantes e integridad.

    \item \textbf{Monitoreo y observabilidad.} Registra métricas, tiempos,
    errores, volumen, retraso y estado de cada etapa.

    \item \textbf{Metadatos, linaje y gobierno.} Documenta el origen, las
    transformaciones, los responsables y el destino de los datos.

    \item \textbf{Seguridad.} Incluye autenticación, autorización, cifrado,
    enmascaramiento, auditoría y protección de datos sensibles.
\end{enumerate}

\section*{4. Estructura de una arquitectura de Data Pipeline}

La estructura general de una arquitectura de Data Pipeline puede representarse
de la siguiente manera:

\begin{center}
\textbf{Fuentes} $\longrightarrow$ \textbf{Ingestión}
$\longrightarrow$ \textbf{Transporte}
$\longrightarrow$ \textbf{Procesamiento}
$\longrightarrow$ \textbf{Almacenamiento o consumo}
\end{center}

\begin{enumerate}
    \item \textbf{Capa de fuentes:} contiene sistemas transaccionales,
    aplicaciones, archivos, API, sensores, registros y servicios externos.

    \item \textbf{Capa de ingestión:} extrae o recibe los datos sin afectar el
    funcionamiento de las fuentes. Puede aplicar cargas completas, cargas
    incrementales o captura de cambios.

    \item \textbf{Capa de transporte:} utiliza colas o plataformas de eventos
    para entregar la información de forma ordenada y resistente a fallas.

    \item \textbf{Capa de procesamiento:} ejecuta reglas de transformación,
    limpieza, integración, agregación y enriquecimiento. Puede trabajar con
    lotes o flujos continuos.

    \item \textbf{Capa de almacenamiento:} guarda datos crudos, intermedios o
    procesados en bases de datos, almacenes analíticos, lagos o Lakehouses.

    \item \textbf{Capa de servicio y consumo:} entrega los resultados a
    tableros, reportes, aplicaciones, API, sistemas de alertas y modelos de
    inteligencia artificial.
\end{enumerate}

La orquestación, seguridad, calidad, metadatos, linaje y observabilidad deben
funcionar de manera transversal. Una arquitectura robusta también incluye una
ruta para los registros que no pueden procesarse, reintentos controlados y
mecanismos para reiniciar el flujo sin perder ni duplicar información.

\section*{5. Características de un Data Pipeline}

Un Data Pipeline moderno debe poseer las siguientes características:

\begin{itemize}
    \item \textbf{Automatización:} minimiza tareas manuales y ejecuta el flujo de manera repetible.
    \item \textbf{Escalabilidad:} aumenta su capacidad ante mayores volúmenes o velocidades de datos.
    \item \textbf{Confiabilidad:} entrega información completa y consistente.
    \item \textbf{Tolerancia a fallas:} utiliza reintentos, puntos de control y recuperación automática.
    \item \textbf{Baja latencia:} procesa los datos con la rapidez requerida por el caso de uso.
    \item \textbf{Idempotencia:} evita resultados distintos cuando una operación debe repetirse.
    \item \textbf{Calidad de datos:} incorpora validación, limpieza y manejo de registros incorrectos.
    \item \textbf{Observabilidad:} proporciona métricas, registros, alertas y trazabilidad.
    \item \textbf{Seguridad:} protege la información durante el transporte, procesamiento y almacenamiento.
    \item \textbf{Modularidad:} divide el flujo en componentes que pueden mantenerse y sustituirse con facilidad.
    \item \textbf{Interoperabilidad:} conecta tecnologías, formatos y plataformas diferentes.
    \item \textbf{Gobierno:} documenta propiedad, linaje, reglas y uso de la información.
\end{itemize}

No todos los pipelines necesitan procesamiento exactamente una vez. La garantía
de entrega debe seleccionarse según el riesgo del caso de uso: una alerta de
fraude o un movimiento financiero requiere controles más estrictos que una
carga estadística que puede reconstruirse.

\section*{6. Ventajas del uso de Data Pipelines}

\begin{itemize}
    \item Automatizan la recolección, movimiento y preparación de los datos.
    \item Reducen errores provocados por actividades manuales.
    \item Entregan información más reciente a usuarios y aplicaciones.
    \item Mejoran la consistencia y calidad de los datos.
    \item Integran fuentes aisladas y disminuyen los silos de información.
    \item Permiten reutilizar procesos de transformación.
    \item Facilitan el crecimiento de las plataformas de datos.
    \item Proporcionan trazabilidad desde el origen hasta el destino.
    \item Aceleran la elaboración de reportes y productos de datos.
    \item Habilitan análisis en tiempo real y decisiones oportunas.
    \item Suministran datos preparados para inteligencia artificial y aprendizaje automático.
    \item Reducen costos operativos cuando se diseñan y monitorean adecuadamente.
\end{itemize}

\section*{7. ¿Por qué son necesarios los Data Pipelines?}

Los Data Pipelines son necesarios porque los datos suelen encontrarse
distribuidos entre sistemas con estructuras, velocidades y formatos diferentes.
Moverlos manualmente resulta lento, propenso a errores y difícil de repetir a
gran escala. Además, la información cruda rara vez está lista para utilizarse
directamente.

Un pipeline establece un proceso controlado para integrar fuentes, aplicar
reglas de calidad y entregar datos a los consumidores adecuados. Esto permite
que las áreas de negocio trabajen con información actualizada, que las
aplicaciones intercambien datos y que los modelos de inteligencia artificial
reciban entradas confiables.

También son necesarios para cumplir niveles de servicio y requisitos de
seguridad. Mediante el monitoreo, linaje y registro de errores, una organización
puede conocer qué ocurrió con cada conjunto de datos y corregir fallas antes de
que afecten decisiones o procesos operativos.

\section*{8. Diferencia entre un Data Pipeline y un proceso ETL}

ETL significa \textit{Extract, Transform, Load}: extraer, transformar y cargar.
Es un patrón específico en el que los datos se extraen de las fuentes, se
transforman antes de llegar al destino y finalmente se cargan. Tradicionalmente
se asocia con cargas por lotes hacia un Data Warehouse, aunque también puede
implementarse en tiempo real.

Un Data Pipeline es un concepto más amplio. Puede incluir un proceso ETL, pero
también puede:

\begin{itemize}
    \item Mover datos sin transformarlos.
    \item Utilizar el patrón ELT y transformar después de la carga.
    \item Procesar eventos continuamente.
    \item Conectar varias fuentes y varios destinos.
    \item Ejecutar algoritmos, validaciones o modelos de aprendizaje automático.
    \item Alimentar aplicaciones operacionales, alertas y servicios, no solamente almacenes de datos.
    \item Incluir orquestación, monitoreo, linaje, seguridad y recuperación ante fallas.
\end{itemize}

Por lo tanto, todo proceso ETL puede considerarse un Data Pipeline, pero no todo
Data Pipeline es un proceso ETL.

\section*{9. Casos de uso comunes de los Data Pipelines}

\begin{itemize}
    \item \textbf{Inteligencia de negocios:} consolidar información en un Data Warehouse para elaborar indicadores y tableros.
    \item \textbf{Data Lakes y Lakehouses:} ingerir datos crudos y producir conjuntos depurados para análisis.
    \item \textbf{Replicación mediante CDC:} mantener sincronizadas bases de datos y plataformas analíticas.
    \item \textbf{Detección de fraude:} evaluar transacciones en tiempo real y generar alertas.
    \item \textbf{Comercio electrónico:} actualizar inventarios, pedidos, precios y recomendaciones.
    \item \textbf{Internet de las cosas:} procesar telemetría de sensores, vehículos y maquinaria.
    \item \textbf{Monitoreo de sistemas:} recolectar registros y métricas para detectar fallas o amenazas.
    \item \textbf{Aprendizaje automático:} preparar datos de entrenamiento, generar variables y alimentar predicciones.
    \item \textbf{Migraciones:} trasladar información entre sistemas locales y plataformas de nube.
    \item \textbf{Vista integral del cliente:} combinar datos de ventas, atención, marketing y navegación.
    \item \textbf{Salud:} integrar expedientes, dispositivos y resultados clínicos bajo controles de privacidad.
    \item \textbf{IA generativa:} actualizar repositorios documentales y bases vectoriales utilizadas por soluciones RAG.
\end{itemize}

\section*{10. Tecnologías o herramientas utilizadas para construir Data Pipelines}

La selección de herramientas depende del volumen, latencia, fuentes, destinos,
presupuesto y conocimientos del equipo. Algunas opciones comunes son:

\begin{itemize}
    \item \textbf{Ingestión y conectores:} Apache NiFi, Airbyte, Fivetran,
    Kafka Connect, Debezium y Striim.

    \item \textbf{Mensajería y eventos:} Apache Kafka, Apache Pulsar,
    RabbitMQ, Amazon Kinesis, Azure Event Hubs y Google Pub/Sub.

    \item \textbf{Procesamiento distribuido:} Apache Spark, Apache Flink,
    Apache Beam y Hazelcast.

    \item \textbf{Orquestación:} Apache Airflow, Dagster, Prefect, Azure Data
    Factory y AWS Glue Workflows.

    \item \textbf{Transformación:} SQL, dbt, Python, Pandas y herramientas
    administradas de integración.

    \item \textbf{Almacenamiento:} PostgreSQL, Amazon S3, Azure Data Lake
    Storage, Google Cloud Storage, Snowflake, BigQuery, Redshift y Databricks.

    \item \textbf{Calidad de datos:} Great Expectations, Soda y Deequ.

    \item \textbf{Catálogo y linaje:} DataHub, OpenMetadata, Apache Atlas y
    OpenLineage.

    \item \textbf{Monitoreo:} Prometheus, Grafana, herramientas de registro y
    servicios de observabilidad de los proveedores de nube.

    \item \textbf{Contenedores e infraestructura:} Docker, Kubernetes y
    herramientas de infraestructura como código.
\end{itemize}

No existe una herramienta universalmente superior. Una arquitectura sencilla,
observable y acorde con las necesidades reales suele ser preferible a combinar
muchas tecnologías sin una justificación clara. Antes de elegir una plataforma
se deben definir la latencia esperada, el volumen, las garantías de entrega, la
seguridad, el costo y la experiencia del equipo.

\section*{Conclusión}

Los Data Pipelines constituyen la infraestructura que conecta la generación de
datos con su utilización. Su valor no se limita a transportar información:
también automatizan transformaciones, aplican controles de calidad, proporcionan
trazabilidad y permiten responder ante fallas. Un pipeline bien diseñado debe
ser escalable, seguro, observable y adecuado a la latencia del caso de uso. La
elección entre lotes, streaming, ETL, ELT o una arquitectura híbrida debe partir
de las necesidades del negocio y no solamente de la disponibilidad de una
tecnología.
